\nonstopmode{}
\documentclass[a4paper]{book}
\usepackage[times,inconsolata,hyper]{Rd}
\usepackage{makeidx}
\makeatletter\@ifl@t@r\fmtversion{2018/04/01}{}{\usepackage[utf8]{inputenc}}\makeatother
% \usepackage{graphicx} % @USE GRAPHICX@
\makeindex{}
\begin{document}
\chapter*{}
\begin{center}
{\textbf{\huge Package `ndm'}}
\par\bigskip{\large \today}
\end{center}
\ifthenelse{\boolean{Rd@use@hyper}}{\hypersetup{pdftitle = {ndm: Neural Disease Modeling Software}}}{}
\ifthenelse{\boolean{Rd@use@hyper}}{\hypersetup{pdfauthor = {Connor T. Jerzak}}}{}
\begin{description}
\raggedright{}
\item[Title]\AsIs{Neural Disease Modeling Software}
\item[Version]\AsIs{0.0.1}
\item[Author]\AsIs{Connor T. Jerzak [aut, cre]}
\item[Maintainer]\AsIs{Connor T. Jerzak }\email{connor.jerzak@austin.utexas.edu}\AsIs{}
\item[Description]\AsIs{Provides an R-first interface for neural disease modeling with a
reticulate-managed JAX backend. The package includes backend bootstrap
helpers, built-in epidemiological model specifications, TFRecord-consuming
data interfaces, package-managed runtime helpers, and package-native
orchestration helpers for model execution.}
\item[Depends]\AsIs{R (>= 4.1.0)}
\item[Imports]\AsIs{data.table,
ndmdatasets,
rlang,
reticulate,
utils}
\item[Suggests]\AsIs{fastmatch,
progress,
rrapply,
testthat (>= 3.0.0),
zip,
zoo}
\item[License]\AsIs{GPL-3}
\item[Encoding]\AsIs{UTF-8}
\item[LazyData]\AsIs{false}
\item[Config/testthat/edition]\AsIs{3}
\item[RoxygenNote]\AsIs{7.3.3}
\end{description}
\Rdcontents{Contents}
\HeaderA{ndm-package}{Neural Disease Modeling Software}{ndm.Rdash.package}
\aliasA{ndm}{ndm-package}{ndm}
%
\begin{Description}
\code{ndm} is an R-first package for neural disease modeling with a
reticulate-managed JAX backend. The package currently provides backend setup,
epidemic model spec import/export, TFRecord readers, package-managed runtime
helpers, and package-native runtime/model helpers.
\end{Description}
%
\begin{Details}
The current implementation preserves both historical model families,
\code{"DecoderOnly"} and \code{"NeuralODE"}, while defaulting to
\code{"DecoderOnly"}. The supported backbone is the transformer path only.
Latent attention and Mamba are intentionally excluded.

The package treats the structured \code{ndm\_model\_spec} object as the canonical
internal representation and supports \code{.tex} import/export for compatibility
with the original epidemic specification workflow.
\end{Details}
%
\begin{Section}{Key functions}
\begin{itemize}

\item{} \code{\LinkA{ndm\_create\_config}{ndm.Rul.create.Rul.config}}
\item{} \code{\LinkA{ndm\_build\_backend}{ndm.Rul.build.Rul.backend}}
\item{} \code{\LinkA{ndm\_initialize\_backend}{ndm.Rul.initialize.Rul.backend}}
\item{} \code{\LinkA{ndm\_model\_spec}{ndm.Rul.model.Rul.spec}}
\item{} \code{\LinkA{ndm\_model\_spec\_from\_tex}{ndm.Rul.model.Rul.spec.Rul.from.Rul.tex}}
\item{} \code{\LinkA{ndm\_model\_spec\_to\_tex}{ndm.Rul.model.Rul.spec.Rul.to.Rul.tex}}
\item{} \code{\LinkA{ndm\_read\_tfrecord\_dataset}{ndm.Rul.read.Rul.tfrecord.Rul.dataset}}
\item{} \code{\LinkA{ndm\_load\_tfrecord\_bundle}{ndm.Rul.load.Rul.tfrecord.Rul.bundle}}
\item{} \code{\LinkA{ndm\_prepare\_runtime}{ndm.Rul.prepare.Rul.runtime}}
\item{} \code{\LinkA{ndm\_prepare\_data}{ndm.Rul.prepare.Rul.data}}
\item{} \code{\LinkA{ndm\_build\_model}{ndm.Rul.build.Rul.model}}
\item{} \code{\LinkA{ndm\_train}{ndm.Rul.train}}
\item{} \code{\LinkA{ndm\_predict}{ndm.Rul.predict}}
\item{} \code{\LinkA{ndm\_loss}{ndm.Rul.loss}}
\item{} \code{\LinkA{ndm\_fit}{ndm.Rul.fit}}

\end{itemize}

\end{Section}
%
\begin{Author}
Connor T. Jerzak \email{connor.jerzak@austin.utexas.edu}
\end{Author}
\HeaderA{ndm\_build\_model}{Build and train models with package-managed runtime stages}{ndm.Rul.build.Rul.model}
\aliasA{ndm\_fit}{ndm\_build\_model}{ndm.Rul.fit}
\aliasA{ndm\_train}{ndm\_build\_model}{ndm.Rul.train}
\aliasA{print.ndm\_model}{ndm\_build\_model}{print.ndm.Rul.model}
\aliasA{print.ndm\_trained\_model}{ndm\_build\_model}{print.ndm.Rul.trained.Rul.model}
%
\begin{Description}
These wrappers layer a small R API over package-managed Phase 1 model build
and training stages.
\end{Description}
%
\begin{Usage}
\begin{verbatim}
ndm_build_model(
  runtime_env,
  model_type = c("DecoderOnly", "NeuralODE"),
  model_spec = NULL,
  backbone = "transformer",
  runtime_globals = list()
)

## S3 method for class 'ndm_model'
print(x, ...)

ndm_train(x, run_define = TRUE, run_loop = TRUE)

## S3 method for class 'ndm_trained_model'
print(x, ...)

ndm_fit(
  config = ndm_create_config(),
  model_spec = NULL,
  data_generator = c("sim", "real", "multidisease"),
  runtime_globals = list(),
  data_globals = list(),
  build_globals = list(),
  train_globals = list(),
  run_define = TRUE,
  run_loop = TRUE
)
\end{verbatim}
\end{Usage}
%
\begin{Arguments}
\begin{ldescription}
\item[\code{runtime\_env}] Runtime environment containing loaded runtime helpers and
data globals.

\item[\code{model\_type}] Model family to build. Either `"DecoderOnly"` or
`"NeuralODE"`.

\item[\code{model\_spec}] Optional `ndm\_model\_spec` object used to override the model
TeX specification supplied to the runtime model builder.

\item[\code{backbone}] Backbone family. Phase 1 supports `"transformer"` only.

\item[\code{runtime\_globals}] Named list of additional globals assigned during the
relevant runtime stage. `ndm\_build\_model()` uses them before building the
model, while `ndm\_fit()` uses them before sourcing the runtime.

\item[\code{x}] Either an `ndm\_model`, an `ndm\_trained\_model`, or a prepared runtime
environment, depending on the function being called.

\item[\code{...}] Unused; included for S3 method compatibility.

\item[\code{run\_define}] Logical scalar indicating whether the training definition
script should be sourced.

\item[\code{run\_loop}] Logical scalar indicating whether the training loop script
should be sourced.

\item[\code{config}] An object of class `ndm\_config`, usually created by
`ndm\_create\_config()`.

\item[\code{data\_generator}] Which local data generator to source before calling
`ndm\_build\_model()` inside `ndm\_fit()`.

\item[\code{data\_globals}] Named list of globals assigned before sourcing the data
generator in `ndm\_fit()`.

\item[\code{build\_globals}] Named list of globals assigned before building the model
in `ndm\_fit()`.

\item[\code{train\_globals}] Named list of globals assigned immediately before
training in `ndm\_fit()`.
\end{ldescription}
\end{Arguments}
%
\begin{Value}
`ndm\_build\_model()` returns an object of class `ndm\_model`.
`ndm\_train()` and `ndm\_fit()` return an object of class
`ndm\_trained\_model`.

`print.ndm\_model()` prints a brief summary of an `ndm\_model` object
and invisibly returns `x`.

`print.ndm\_trained\_model()` prints a brief summary of an
`ndm\_trained\_model` object and invisibly returns `x`.
\end{Value}
%
\begin{Examples}
\begin{ExampleCode}
cfg <- ndm_create_config()
spec <- ndm_model_spec()
## Not run: 
runtime_env <- ndm_prepare_runtime(cfg)
ndm_prepare_data(runtime_env, generator = "sim")
model <- ndm_build_model(runtime_env, model_spec = spec)
trained <- ndm_train(model)

## End(Not run)

\end{ExampleCode}
\end{Examples}
\HeaderA{ndm\_check\_backend}{Provision and initialize the Python backend}{ndm.Rul.check.Rul.backend}
\aliasA{ndm\_backend\_modules}{ndm\_check\_backend}{ndm.Rul.backend.Rul.modules}
\aliasA{ndm\_build\_backend}{ndm\_check\_backend}{ndm.Rul.build.Rul.backend}
\aliasA{ndm\_initialize\_backend}{ndm\_check\_backend}{ndm.Rul.initialize.Rul.backend}
%
\begin{Description}
These helpers manage the reticulate-backed Python environment used by
ndm. Use `ndm\_build\_backend()` to provision dependencies,
`ndm\_check\_backend()` to verify availability, and
`ndm\_initialize\_backend()` to import the active JAX stack into R.
\end{Description}
%
\begin{Usage}
\begin{verbatim}
ndm_check_backend(
  conda_env = "ndm_software_env",
  conda = "auto",
  modules = c("jax", "numpy", "optax", "equinox", "diffrax")
)

ndm_build_backend(
  conda_env = "ndm_software_env",
  conda = "auto",
  python_version = "3.13",
  include_tensorflow = TRUE,
  extra_packages = c("optax", "equinox", "diffrax")
)

ndm_initialize_backend(
  conda_env = "ndm_software_env",
  conda = "auto",
  conda_env_required = TRUE,
  float_type = c("32", "64"),
  import_tensorflow = TRUE
)

ndm_backend_modules()
\end{verbatim}
\end{Usage}
%
\begin{Arguments}
\begin{ldescription}
\item[\code{conda\_env}] Name of the conda environment that should contain the JAX
runtime.

\item[\code{conda}] Conda executable to use. `"auto"` delegates discovery to
reticulate.

\item[\code{modules}] Character vector of Python module names that must be importable
for the backend to be considered healthy.

\item[\code{python\_version}] Python version requested when creating the conda
environment.

\item[\code{include\_tensorflow}] Logical scalar indicating whether TensorFlow should
be installed alongside JAX.

\item[\code{extra\_packages}] Additional Python packages to install with `pip`.

\item[\code{conda\_env\_required}] Passed through to `reticulate::use\_condaenv()` as
the `required` flag.

\item[\code{float\_type}] Floating point precision used when configuring JAX. Either
`"32"` or `"64"`.

\item[\code{import\_tensorflow}] Logical scalar indicating whether TensorFlow should
be imported when it is available in the selected environment.
\end{ldescription}
\end{Arguments}
%
\begin{Value}
`ndm\_check\_backend()` returns `TRUE` invisibly when the requested
environment and modules are available. It returns `NULL` after printing a
diagnostic message when the backend is not ready.

`ndm\_build\_backend()` invisibly returns the conda environment name
after attempting to provision the requested packages.

`ndm\_initialize\_backend()` returns an object of class
`ndm\_backend`. The returned list contains imported Python modules, the
configured float type, and a small compatibility shim used by the legacy
runtime.

`ndm\_backend\_modules()` returns the currently initialized
`ndm\_backend` object.
\end{Value}
%
\begin{Examples}
\begin{ExampleCode}
ndm_check_backend()

## Not run: 
ndm_build_backend()

## End(Not run)

## Not run: 
backend <- ndm_initialize_backend()
backend$default_backend

## End(Not run)

## Not run: 
ndm_initialize_backend()
ndm_backend_modules()

## End(Not run)

\end{ExampleCode}
\end{Examples}
\HeaderA{ndm\_create\_config}{Create runtime configuration objects}{ndm.Rul.create.Rul.config}
\aliasA{print.ndm\_config}{ndm\_create\_config}{print.ndm.Rul.config}
%
\begin{Description}
Configuration objects collect the runtime options that are threaded through
the higher-level orchestration helpers such as `ndm\_prepare\_runtime()` and
`ndm\_fit()`.
\end{Description}
%
\begin{Usage}
\begin{verbatim}
ndm_create_config(
  model_type = c("DecoderOnly", "NeuralODE"),
  backbone = "transformer",
  float_type = c("32", "64"),
  force_to_gpu = TRUE,
  resave_tfrecords = FALSE,
  gpu_mem_frac = NULL,
  ...
)

## S3 method for class 'ndm_config'
print(x, ...)
\end{verbatim}
\end{Usage}
%
\begin{Arguments}
\begin{ldescription}
\item[\code{model\_type}] Model family metadata. Use `"DecoderOnly"` or
`"NeuralODE"`.

\item[\code{backbone}] Backbone family. Phase 1 supports `"transformer"` only.

\item[\code{float\_type}] Floating point precision used when initializing the backend.
Use `"32"` or `"64"`.

\item[\code{force\_to\_gpu}] Logical scalar indicating whether the runtime should try
to place arrays on GPU when available.

\item[\code{resave\_tfrecords}] Logical scalar preserved for compatibility with
existing run workflows. Multidisease workflows reject `TRUE` because the
legacy TFRecord-regeneration path has been retired.

\item[\code{gpu\_mem\_frac}] Optional GPU memory fraction forwarded into the runtime
bootstrap code.

\item[\code{...}] Unused; included for S3 method compatibility.

\item[\code{x}] An `ndm\_config` object.
\end{ldescription}
\end{Arguments}
%
\begin{Value}
`ndm\_create\_config()` returns an object of class `ndm\_config`.

`print.ndm\_config()` prints a brief summary of an `ndm\_config`
object and invisibly returns `x`.
\end{Value}
%
\begin{Examples}
\begin{ExampleCode}
cfg <- ndm_create_config()
print(cfg)

\end{ExampleCode}
\end{Examples}
\HeaderA{ndm\_create\_real\_run\_config}{Package-native run configurations and workflow runners}{ndm.Rul.create.Rul.real.Rul.run.Rul.config}
\aliasA{ndm\_create\_multidisease\_run\_config}{ndm\_create\_real\_run\_config}{ndm.Rul.create.Rul.multidisease.Rul.run.Rul.config}
\aliasA{ndm\_create\_sim\_run\_config}{ndm\_create\_real\_run\_config}{ndm.Rul.create.Rul.sim.Rul.run.Rul.config}
\aliasA{ndm\_run\_multidisease}{ndm\_create\_real\_run\_config}{ndm.Rul.run.Rul.multidisease}
\aliasA{ndm\_run\_real}{ndm\_create\_real\_run\_config}{ndm.Rul.run.Rul.real}
\aliasA{ndm\_run\_sim}{ndm\_create\_real\_run\_config}{ndm.Rul.run.Rul.sim}
%
\begin{Description}
Package-native configuration constructors and workflow runners for the
self-contained real, simulation, and multidisease execution paths.
\end{Description}
%
\begin{Usage}
\begin{verbatim}
ndm_create_real_run_config(
  project_root = getwd(),
  analysis_name = "RealApril15",
  grid = NULL,
  grid_file = file.path("Data", "RunGrids", "RealGrids", sprintf("RealGrid_%s.csv",
    analysis_name)),
  outer = 3L,
  model_type = NULL,
  respect_grid_model_type = FALSE,
  resave_tfrecords = FALSE,
  tfrecord_dir = file.path("Data", "RunTFRecords", "RealTFRecords",
    analysis_name),
  raw_data_dir = file.path("Data", "MainData"),
  outcome_metric = "inc_death",
  data_subset = "high_income",
  dry_run = FALSE
)

ndm_create_sim_run_config(
  project_root = getwd(),
  analysis_name = "BigSimsLatest",
  grid = NULL,
  grid_file = file.path("Data", "RunGrids", "SimGrids", sprintf("SimGrid_%s.csv",
    analysis_name)),
  outer = 1L,
  model_type = NULL,
  respect_grid_model_type = FALSE,
  resave_tfrecords = TRUE,
  tfrecord_dir = file.path("Data", "RunTFRecords", "SimTFRecords",
    analysis_name),
  dry_run = FALSE
)

ndm_create_multidisease_run_config(
  project_root = getwd(),
  analysis_name = "RealLatest",
  grid = NULL,
  grid_file = file.path("Data", "RunGrids", "RealGrids", sprintf("RealGrid_%s.csv",
    analysis_name)),
  outer = 1L,
  model_type = NULL,
  respect_grid_model_type = FALSE,
  resave_tfrecords = FALSE,
  tfrecord_dir = file.path("Data", "RunTFRecords", "RealTFRecords",
    analysis_name),
  outcome_metric = "CountValue",
  data_subset = "all",
  disease_names = c("Covid", "Flu"),
  data_format = "IHME",
  dry_run = FALSE
)

ndm_run_real(config = ndm_create_real_run_config())

ndm_run_sim(config = ndm_create_sim_run_config())

ndm_run_multidisease(config = ndm_create_multidisease_run_config())
\end{verbatim}
\end{Usage}
%
\begin{Arguments}
\begin{ldescription}
\item[\code{project\_root}] Working project directory used for grids, TFRecords, and outputs.

\item[\code{analysis\_name}] Analysis label threaded through output paths.

\item[\code{grid}] Optional in-memory grid data frame. When supplied, \code{grid\_file} must be \code{NULL}.

\item[\code{grid\_file}] Optional CSV path for the run grid.

\item[\code{outer}] Integer vector of outer-iteration rows to execute.

\item[\code{model\_type}] Optional model family override.

\item[\code{respect\_grid\_model\_type}] Logical scalar indicating whether grid rows may override the requested model type.

\item[\code{resave\_tfrecords}] Logical scalar controlling TFRecord regeneration.
Supported for real and simulation workflows. Multidisease workflows reject
\code{TRUE} because the legacy regeneration path has been retired.

\item[\code{tfrecord\_dir}] Optional TFRecord output directory.

\item[\code{raw\_data\_dir}] Real-data input directory.

\item[\code{outcome\_metric}] Outcome metric name for real-data or multidisease runs.

\item[\code{data\_subset}] Optional real-data subset name.

\item[\code{disease\_names}] Character vector of multidisease names.

\item[\code{data\_format}] Multidisease data format.

\item[\code{dry\_run}] Logical scalar indicating whether the run should stop after resolving paths and grid rows.

\item[\code{config}] A package-native run configuration created by the matching \code{ndm\_create\_*\_run\_config()} helper.
\end{ldescription}
\end{Arguments}
%
\begin{Value}
The \code{ndm\_create\_*\_run\_config()} helpers return classed lists describing
the requested workflow. The \code{ndm\_run\_*()} functions return the underlying
workflow result, or a dry-run preview when \code{config\$dry\_run} is \code{TRUE}.
\end{Value}
%
\begin{Examples}
\begin{ExampleCode}
## Not run: 
cfg <- ndm_create_sim_run_config(
  project_root = getwd(),
  grid_file = "Data/RunGrids/SimGrids/SimGrid_BigSimsLatest.csv",
  outer = 1L,
  dry_run = TRUE
)

ndm_run_sim(cfg)

## End(Not run)
\end{ExampleCode}
\end{Examples}
\HeaderA{ndm\_create\_tfrecord\_parser}{Create and read TensorFlow TFRecord datasets}{ndm.Rul.create.Rul.tfrecord.Rul.parser}
\aliasA{ndm\_collect\_tfrecord\_batches}{ndm\_create\_tfrecord\_parser}{ndm.Rul.collect.Rul.tfrecord.Rul.batches}
\aliasA{ndm\_read\_tfrecord\_dataset}{ndm\_create\_tfrecord\_parser}{ndm.Rul.read.Rul.tfrecord.Rul.dataset}
%
\begin{Description}
These wrappers construct parsers for the ndm TFRecord schema and use
them to build TensorFlow datasets or eagerly collect batches into R.
\end{Description}
%
\begin{Usage}
\begin{verbatim}
ndm_create_tfrecord_parser(field_names, dtype_map = NULL, tensorflow = NULL)

ndm_read_tfrecord_dataset(
  file,
  batch_size,
  field_names,
  dtype_map = NULL,
  max_examples = NULL,
  shuffle = FALSE,
  buffer_scaler = 10L,
  tensorflow = NULL
)

ndm_collect_tfrecord_batches(
  file,
  batch_size,
  field_names,
  dtype_map = NULL,
  max_batches = NULL,
  max_examples = NULL,
  shuffle = FALSE,
  buffer_scaler = 10L,
  tensorflow = NULL
)
\end{verbatim}
\end{Usage}
%
\begin{Arguments}
\begin{ldescription}
\item[\code{field\_names}] Character vector of TFRecord field names that should be
parsed from each example.

\item[\code{dtype\_map}] Optional named vector or list mapping `field\_names` to
TensorFlow dtype names or TensorFlow dtype objects.

\item[\code{tensorflow}] Optional TensorFlow module. When omitted, ndm uses
the active backend's TensorFlow module when available and otherwise
imports TensorFlow from the active reticulate environment at call time.

\item[\code{file}] Path to the `.tfrecord` file to read.

\item[\code{batch\_size}] Batch size used when constructing the TensorFlow dataset.

\item[\code{max\_examples}] Optional maximum number of examples to consume from the
file before batching.

\item[\code{shuffle}] Logical scalar indicating whether the dataset should be
shuffled.

\item[\code{buffer\_scaler}] Integer multiplier used to derive the TensorFlow shuffle
buffer size from `batch\_size`.

\item[\code{max\_batches}] Optional maximum number of batches to collect when using
`ndm\_collect\_tfrecord\_batches()`.
\end{ldescription}
\end{Arguments}
%
\begin{Value}
`ndm\_create\_tfrecord\_parser()` returns a parser function compatible
with `tf\$data\$Dataset\$map()`. `ndm\_read\_tfrecord\_dataset()` returns a
TensorFlow dataset object. `ndm\_collect\_tfrecord\_batches()` returns a list
of parsed batches.
\end{Value}
%
\begin{Examples}
\begin{ExampleCode}
fields <- ndm_tfrecord_fields_real()
parser <- ndm_create_tfrecord_parser(fields, ndm_tfrecord_dtype_map("real"))
is.function(parser)

\end{ExampleCode}
\end{Examples}
\HeaderA{ndm\_model\_spec\_presets}{Inspect and create epidemic model specifications}{ndm.Rul.model.Rul.spec.Rul.presets}
\aliasA{ndm\_model\_spec}{ndm\_model\_spec\_presets}{ndm.Rul.model.Rul.spec}
\aliasA{ndm\_model\_spec\_from\_tex}{ndm\_model\_spec\_presets}{ndm.Rul.model.Rul.spec.Rul.from.Rul.tex}
\aliasA{ndm\_model\_spec\_to\_tex}{ndm\_model\_spec\_presets}{ndm.Rul.model.Rul.spec.Rul.to.Rul.tex}
\aliasA{print.ndm\_model\_spec}{ndm\_model\_spec\_presets}{print.ndm.Rul.model.Rul.spec}
%
\begin{Description}
The package ships a small registry of built-in compartmental model
specifications and also supports importing custom `.tex` model definitions
from disk. The structured `ndm\_model\_spec` object is the canonical
representation used by the runtime wrappers.
\end{Description}
%
\begin{Usage}
\begin{verbatim}
ndm_model_spec_presets()

ndm_model_spec(
  preset = NULL,
  compartments = c("seir", "seirs"),
  dynamic_beta = NULL,
  dynamic_global = FALSE,
  multi_outcome = FALSE,
  model_type = c("DecoderOnly", "NeuralODE"),
  time_varying_terms = NULL,
  endogenous_terms = NULL,
  tex_path = NULL
)

ndm_model_spec_from_tex(path, model_type = c("DecoderOnly", "NeuralODE"))

ndm_model_spec_to_tex(spec, path = NULL)

## S3 method for class 'ndm_model_spec'
print(x, ...)
\end{verbatim}
\end{Usage}
%
\begin{Arguments}
\begin{ldescription}
\item[\code{preset}] Optional preset name from `ndm\_model\_spec\_presets()`. When
omitted, the remaining arguments are used to infer a built-in preset.

\item[\code{compartments}] Compartment family used when inferring a built-in preset.

\item[\code{dynamic\_beta}] Logical scalar controlling whether the local transmission
rate is time-varying for built-in presets.

\item[\code{dynamic\_global}] Logical scalar controlling whether the global rates are
time-varying for built-in presets.

\item[\code{multi\_outcome}] Logical scalar indicating whether the observation model
should include multiple outcomes.

\item[\code{model\_type}] Model family metadata attached to the returned
specification. Either `"DecoderOnly"` or `"NeuralODE"`.

\item[\code{time\_varying\_terms}] Optional character vector overriding the default
time-varying term metadata stored on the returned spec.

\item[\code{endogenous\_terms}] Optional character vector overriding the default
endogenous term metadata stored on the returned spec.

\item[\code{tex\_path}] Optional path to a custom `.tex` file. When supplied,
`ndm\_model\_spec()` delegates to `ndm\_model\_spec\_from\_tex()`.

\item[\code{path}] Path to a `.tex` model specification file on disk.

\item[\code{spec}] An `ndm\_model\_spec` object.

\item[\code{x}] An `ndm\_model\_spec` object.

\item[\code{...}] Unused; included for S3 method compatibility.
\end{ldescription}
\end{Arguments}
%
\begin{Value}
`ndm\_model\_spec\_presets()` returns a data frame describing the
built-in presets bundled with the package.

`ndm\_model\_spec()` returns an object of class `ndm\_model\_spec`.
Built-in presets include metadata such as the packaged source path and the
normalized TeX contents.

`ndm\_model\_spec\_from\_tex()` returns an `ndm\_model\_spec` object whose
`source\_path` points at the supplied file. When the basename matches a
built-in preset, the preset metadata is reused.

`ndm\_model\_spec\_to\_tex()` returns the TeX text as a single character
string when `path` is `NULL`. When `path` is supplied, it writes the
specification to disk and invisibly returns the normalized output path.

`print.ndm\_model\_spec()` prints a brief summary of an
`ndm\_model\_spec` object and invisibly returns `x`.
\end{Value}
%
\begin{Examples}
\begin{ExampleCode}
ndm_model_spec_presets()

spec <- ndm_model_spec()
print(spec)

spec <- ndm_model_spec()
tex_path <- tempfile(fileext = ".tex")
ndm_model_spec_to_tex(spec, tex_path)
ndm_model_spec_from_tex(tex_path)

spec <- ndm_model_spec()
tex <- ndm_model_spec_to_tex(spec)
nchar(tex) > 0

\end{ExampleCode}
\end{Examples}
\HeaderA{ndm\_new\_runtime\_env}{Create and populate runtime environments}{ndm.Rul.new.Rul.runtime.Rul.env}
\aliasA{ndm\_set\_runtime\_globals}{ndm\_new\_runtime\_env}{ndm.Rul.set.Rul.runtime.Rul.globals}
%
\begin{Description}
These helpers manage isolated environments used to load package-managed
runtime code and seed it with R values.
\end{Description}
%
\begin{Usage}
\begin{verbatim}
ndm_new_runtime_env(parent = globalenv())

ndm_set_runtime_globals(env, values, overwrite = TRUE)
\end{verbatim}
\end{Usage}
%
\begin{Arguments}
\begin{ldescription}
\item[\code{parent}] Parent environment for a new runtime environment. Defaults to
`globalenv()` so package-managed runtime code can resolve attached-package
functions consistently.

\item[\code{env}] Runtime environment that should receive new bindings.

\item[\code{values}] Named list of values to assign into `env`.

\item[\code{overwrite}] Logical scalar indicating whether existing bindings in `env`
should be overwritten.
\end{ldescription}
\end{Arguments}
%
\begin{Value}
`ndm\_new\_runtime\_env()` returns an environment of class
`ndm\_runtime\_env`. `ndm\_set\_runtime\_globals()` invisibly returns `env`.
\end{Value}
%
\begin{Examples}
\begin{ExampleCode}
env <- ndm_new_runtime_env()
ndm_set_runtime_globals(env, list(example_value = 1L))
env$example_value

\end{ExampleCode}
\end{Examples}
\HeaderA{ndm\_predict}{Generate predictions or evaluate the training loss}{ndm.Rul.predict}
\aliasA{ndm\_loss}{ndm\_predict}{ndm.Rul.loss}
%
\begin{Description}
These helpers call the runtime prediction and loss functions stored in a
prepared runtime environment.
\end{Description}
%
\begin{Usage}
\begin{verbatim}
ndm_predict(x, batch = NULL, inference = TRUE, seed = 1L, update_state = TRUE)

ndm_loss(
  x,
  batch = NULL,
  y = NULL,
  y_mask = NULL,
  iteration = 1L,
  seed = 1L,
  update_state = TRUE
)
\end{verbatim}
\end{Usage}
%
\begin{Arguments}
\begin{ldescription}
\item[\code{x}] Either an `ndm\_model`, an `ndm\_trained\_model`, or a prepared runtime
environment that already contains the required runtime objects.

\item[\code{batch}] Optional batch object. Supply either a named TFRecord-style list
or the four-element packaged list returned by `ndm\_batch\_to\_model\_inputs()`.
When `NULL`, the runtime must already contain `batch\_l\_cal`.

\item[\code{inference}] Logical scalar indicating whether `ndm\_predict()` should use
the inference prediction path or the training prediction path.

\item[\code{seed}] Integer seed used to derive the per-example JAX RNG keys.

\item[\code{update\_state}] Logical scalar indicating whether the runtime `state`
object should be updated with the returned state.

\item[\code{y}] Optional observed outcome tensor passed to `ndm\_loss()`. When
omitted, `YTrue\_out` is inferred from `batch`.

\item[\code{y\_mask}] Optional outcome mask tensor passed to `ndm\_loss()`. When
omitted, `YTrue\_out\_mask` is inferred from `batch`.

\item[\code{iteration}] Training iteration number forwarded to the runtime loss
function.
\end{ldescription}
\end{Arguments}
%
\begin{Value}
`ndm\_predict()` returns the prediction object produced by the
runtime. `ndm\_loss()` returns the loss object produced by the runtime.
\end{Value}
%
\begin{Examples}
\begin{ExampleCode}
## Not run: 
preds <- ndm_predict(model, batch = batch_l_cal)
loss <- ndm_loss(model, batch = batch_l_cal)

## End(Not run)

\end{ExampleCode}
\end{Examples}
\HeaderA{ndm\_prepare\_runtime}{Prepare a local runtime for execution}{ndm.Rul.prepare.Rul.runtime}
\aliasA{ndm\_prepare\_data}{ndm\_prepare\_runtime}{ndm.Rul.prepare.Rul.data}
%
\begin{Description}
These wrappers load package-owned runtime components into an isolated
environment and then source the selected data generator.
\end{Description}
%
\begin{Usage}
\begin{verbatim}
ndm_prepare_runtime(
  config = ndm_create_config(),
  runtime_env = ndm_new_runtime_env(),
  runtime_globals = list()
)

ndm_prepare_data(
  runtime_env,
  generator = c("sim", "real", "multidisease"),
  runtime_globals = list()
)
\end{verbatim}
\end{Usage}
%
\begin{Arguments}
\begin{ldescription}
\item[\code{config}] An object of class `ndm\_config`, usually created by
`ndm\_create\_config()`.

\item[\code{runtime\_env}] Runtime environment that should receive the sourced
objects.

\item[\code{runtime\_globals}] Named list of additional bindings to assign before the
runtime code is sourced.

\item[\code{generator}] Which data generator to source into `runtime\_env`.
\end{ldescription}
\end{Arguments}
%
\begin{Value}
`ndm\_prepare\_runtime()` invisibly returns `runtime\_env` after
loading the package-owned helper and backend code.

`ndm\_prepare\_data()` invisibly returns `runtime\_env` after sourcing
the requested data generator.
\end{Value}
%
\begin{Examples}
\begin{ExampleCode}
env <- ndm_new_runtime_env()
cfg <- ndm_create_config()
## Not run: 
ndm_prepare_runtime(cfg, env)

## End(Not run)

env <- ndm_new_runtime_env()
ndm_set_runtime_globals(env, list(example_value = 1L))

\end{ExampleCode}
\end{Examples}
\HeaderA{ndm\_print}{Print a timestamped diagnostic message}{ndm.Rul.print}
%
\begin{Description}
`ndm\_print()` is a small logging helper used throughout the package-facing
orchestration code.
\end{Description}
%
\begin{Usage}
\begin{verbatim}
ndm_print(text, quiet = FALSE)
\end{verbatim}
\end{Usage}
%
\begin{Arguments}
\begin{ldescription}
\item[\code{text}] Text to print.

\item[\code{quiet}] Logical scalar indicating whether output should be suppressed.
\end{ldescription}
\end{Arguments}
%
\begin{Value}
The input `text`, invisibly.
\end{Value}
%
\begin{Examples}
\begin{ExampleCode}
ndm_print("starting")

\end{ExampleCode}
\end{Examples}
\HeaderA{ndm\_save\_model}{Save and restore model artifacts}{ndm.Rul.save.Rul.model}
\aliasA{ndm\_load\_model}{ndm\_save\_model}{ndm.Rul.load.Rul.model}
\aliasA{ndm\_resume\_training}{ndm\_save\_model}{ndm.Rul.resume.Rul.training}
%
\begin{Description}
These helpers persist the public `ndm\_model` / `ndm\_trained\_model` objects to
a versioned artifact directory, restore them into a fresh runtime
environment, and resume training from the restored optimizer state.
\end{Description}
%
\begin{Usage}
\begin{verbatim}
ndm_save_model(x, path, bundle = NULL, overwrite = FALSE)

ndm_load_model(
  path,
  bundle = NULL,
  calibration_source = c("inference", "train")
)

ndm_resume_training(
  path,
  bundle = NULL,
  calibration_source = c("train", "inference"),
  train_globals = list()
)
\end{verbatim}
\end{Usage}
%
\begin{Arguments}
\begin{ldescription}
\item[\code{x}] An object of class `ndm\_model` or `ndm\_trained\_model`.

\item[\code{path}] Artifact directory path.

\item[\code{bundle}] Optional `ndm\_tfrecord\_bundle` used to record TFRecord
references for later resume workflows.

\item[\code{overwrite}] Logical scalar indicating whether an existing artifact
directory should be replaced.

\item[\code{calibration\_source}] Which dataset should supply the preview batch used
to seed runtime globals such as `batch\_l\_cal` when `bundle` is attached.

\item[\code{train\_globals}] Optional named list of runtime globals applied
immediately before `ndm\_resume\_training()` continues the training loop.
\end{ldescription}
\end{Arguments}
%
\begin{Value}
`ndm\_save\_model()` returns the normalized artifact directory path,
invisibly. `ndm\_load\_model()` returns an `ndm\_model` or
`ndm\_trained\_model` matching the saved artifact. `ndm\_resume\_training()`
returns an `ndm\_trained\_model`.
\end{Value}
%
\begin{Examples}
\begin{ExampleCode}
## Not run: 
artifact_dir <- ndm_save_model(trained_model, tempfile("ndm-artifact-"))
restored <- ndm_load_model(artifact_dir)
resumed <- ndm_resume_training(artifact_dir, bundle = tf_bundle)

## End(Not run)

\end{ExampleCode}
\end{Examples}
\HeaderA{ndm\_tfrecord\_fields\_real}{Inspect TFRecord field contracts}{ndm.Rul.tfrecord.Rul.fields.Rul.real}
\aliasA{ndm\_tfrecord\_dtype\_map}{ndm\_tfrecord\_fields\_real}{ndm.Rul.tfrecord.Rul.dtype.Rul.map}
\aliasA{ndm\_tfrecord\_fields\_sim}{ndm\_tfrecord\_fields\_real}{ndm.Rul.tfrecord.Rul.fields.Rul.sim}
%
\begin{Description}
These helpers describe the field names and default TensorFlow dtypes expected
by the packaged TFRecord readers.
\end{Description}
%
\begin{Usage}
\begin{verbatim}
ndm_tfrecord_fields_real()

ndm_tfrecord_fields_sim()

ndm_tfrecord_dtype_map(kind = c("real", "sim"))
\end{verbatim}
\end{Usage}
%
\begin{Arguments}
\begin{ldescription}
\item[\code{kind}] TFRecord schema to inspect. Use `"real"` for observed-data
records and `"sim"` for simulation records.
\end{ldescription}
\end{Arguments}
%
\begin{Value}
`ndm\_tfrecord\_fields\_real()` and `ndm\_tfrecord\_fields\_sim()` return
character vectors of field names. `ndm\_tfrecord\_dtype\_map()` returns a
named character vector mapping each field to its default TensorFlow dtype.
\end{Value}
%
\begin{Examples}
\begin{ExampleCode}
ndm_tfrecord_fields_real()
ndm_tfrecord_dtype_map("sim")

\end{ExampleCode}
\end{Examples}
\HeaderA{ndm\_tf\_batch\_to\_r}{Convert and bundle TFRecord batches}{ndm.Rul.tf.Rul.batch.Rul.to.Rul.r}
\aliasA{ndm\_attach\_tfrecord\_bundle}{ndm\_tf\_batch\_to\_r}{ndm.Rul.attach.Rul.tfrecord.Rul.bundle}
\aliasA{ndm\_batch\_to\_model\_inputs}{ndm\_tf\_batch\_to\_r}{ndm.Rul.batch.Rul.to.Rul.model.Rul.inputs}
\aliasA{ndm\_load\_tfrecord\_bundle}{ndm\_tf\_batch\_to\_r}{ndm.Rul.load.Rul.tfrecord.Rul.bundle}
\aliasA{ndm\_tf\_batch\_to\_jax}{ndm\_tf\_batch\_to\_r}{ndm.Rul.tf.Rul.batch.Rul.to.Rul.jax}
%
\begin{Description}
These helpers convert TensorFlow batches to native R or JAX objects, package
them into the shape expected by the package-managed runtime, and attach
dataset handles to a runtime environment.
\end{Description}
%
\begin{Usage}
\begin{verbatim}
ndm_tf_batch_to_r(batch)

ndm_tf_batch_to_jax(batch, backend = NULL)

ndm_batch_to_model_inputs(batch)

ndm_load_tfrecord_bundle(
  train_file,
  inference_file = NULL,
  batch_size,
  kind = c("real", "sim"),
  dtype_map = NULL,
  tensorflow = NULL,
  max_train_examples = NULL,
  max_inference_examples = NULL,
  shuffle_train = TRUE,
  shuffle_inference = FALSE
)

ndm_attach_tfrecord_bundle(
  env,
  bundle,
  backend = NULL,
  calibration_source = c("inference", "train")
)
\end{verbatim}
\end{Usage}
%
\begin{Arguments}
\begin{ldescription}
\item[\code{batch}] A parsed TFRecord batch or a nested list containing TensorFlow
tensors.

\item[\code{backend}] Optional `ndm\_backend` object. When omitted,
`ndm\_tf\_batch\_to\_jax()` and `ndm\_attach\_tfrecord\_bundle()` use the active
backend from `ndm\_backend\_modules()` when available.

\item[\code{train\_file}] Path to the training TFRecord file.

\item[\code{inference\_file}] Optional path to an inference TFRecord file.

\item[\code{batch\_size}] Batch size used when constructing the training and
inference datasets.

\item[\code{kind}] TFRecord schema to use when reading `train\_file` and
`inference\_file`.

\item[\code{dtype\_map}] Optional named dtype mapping passed through to the TFRecord
readers.

\item[\code{tensorflow}] Optional TensorFlow module. When omitted, ndm uses the
active backend's TensorFlow module when available and otherwise imports
TensorFlow from the active reticulate environment at call time.

\item[\code{max\_train\_examples}] Optional cap on the number of training examples to
read.

\item[\code{max\_inference\_examples}] Optional cap on the number of inference
examples to read.

\item[\code{shuffle\_train}] Logical scalar indicating whether the training dataset
should be shuffled.

\item[\code{shuffle\_inference}] Logical scalar indicating whether the inference
dataset should be shuffled.

\item[\code{env}] Runtime environment that should receive dataset handles and the
preview calibration batch.

\item[\code{bundle}] An object of class `ndm\_tfrecord\_bundle` returned by
`ndm\_load\_tfrecord\_bundle()`.

\item[\code{calibration\_source}] Which dataset should supply the preview batch used
to seed runtime globals such as `batch\_l\_cal`.
\end{ldescription}
\end{Arguments}
%
\begin{Value}
`ndm\_tf\_batch\_to\_r()` recursively converts a TensorFlow batch to R
arrays and lists. `ndm\_tf\_batch\_to\_jax()` recursively converts tensors into
JAX arrays. `ndm\_batch\_to\_model\_inputs()` returns the packaged four-element
list expected by the package-managed model stages. `ndm\_load\_tfrecord\_bundle()`
returns an `ndm\_tfrecord\_bundle` object. `ndm\_attach\_tfrecord\_bundle()`
invisibly returns `env`.
\end{Value}
%
\begin{Examples}
\begin{ExampleCode}
batch <- list(
  XPred = array(0, c(2, 3, 4)),
  XPred_mask = array(TRUE, c(2, 3, 4)),
  Context = array(0, c(2, 1, 4)),
  Context_mask = array(TRUE, c(2, 1, 4)),
  location_id_numeric = matrix(1L, nrow = 2),
  time_id_numeric = matrix(1L, nrow = 2)
)
packaged <- ndm_batch_to_model_inputs(batch)
length(packaged)

\end{ExampleCode}
\end{Examples}
\printindex{}
\end{document}
